\section{Diseño arquitectural}

\subsection{Definición de la arquitectura de software}

La solución que se ha implementado cuenta con un patrón cliente web + BFF + servicios internos. El frontend, desarrollado en React consume rutas del backend FastAPI. El backend actua como BFF se encarga de la extracción del LLM, la fusión de los datos dentro de la sesión, el cálculo de opciones, la persistencia y las respuestas para la interfaz de usuario (canal). Para la creación de la aplicación el archivo \texttt{main.py} se encarga del CORS; \texttt{routes.py} concentra las interfaces; \texttt{session\_service.py} se encarga de las sesiones; \texttt{decision\_engine.py} se encarga de la toma de decisión sobre la mejor sesión; \texttt{json\_repository.py} se encarga de la persistencia de los datos.

Además, se considera una capa de integración para canales de comunicación. En el estado actual del proyecto, el sistema cuenta con un canal simulado para las conversasiones grupales. Para el desarrollo del MVP se contempla la implementación de al menos una integración real, priorizando Discord debido a su simplicidad de integración. La capa quedara diseñada para permitir futuras integraciones ya sea con WhatsApp, Slack u otras plataformas.
Finalmente, se incorporará un servicio de calendario que permitira consolidar e integrar la decisión confirmada en un evento calendarizable, inicialmente en un formato .ics.

El flujo general quedaria asi: mensaje del usuario $\rightarrow$ canal web o externo $\rightarrow$ BFF $\rightarrow$ servicio LLM extractor $\rightarrow$ actualización de sesión $\rightarrow$ motor de decisión $\rightarrow$ respuesta al usuario $\rightarrow$ confirmación $\rightarrow$ generación de evento de calendario.

\begin{figure}[htbp]
    \centering
    \includegraphics[width=0.9\textwidth]{Arquitectura_General.png}
    \caption{Diagrama Arquitectura general}
    \label{fig:arquitectura_general}
\end{figure}

\subsection{Identificación de la integración del LLM en la solución}
El LLM se integra como un componente dentro del sercio \textit{LLM Service}. Es importante que la función no es tomar una decisión final, sino transformar lo mensajes en lenguaje natural a datos estructurados en un formato especifico para que puedan ser utilizados el resto del sistema.

Por ejemplo, si llegase un mensaje como ``Nicolás puede el lunes en la tarde y Matías puede desde las 16'', el LLM debe extraer participantes, días y horarios disponibles. Esta información es devuelta en un formato estructurado en JSON, para que el backend lo procese y lo fusione con el resto de la información de la sesión.

La decisión final depende del servicio \textit{Decision Engine}, el servicio se encarga de calcular de manera deterministica la mejor opcion evaluando la cobertura de los participantes y generando una explicación para cada alternativa generada.

La integración considera mecanismos basicos de seguridad como prompts restrictivos, validación de la salida, correccion manual de los datos reconocidos por el sistema, uso de fallback en caso de falllas del modelo y la confirmación explicita del usuario a la hora de la toma  de decisión final.

\begin{figure}[htbp]
    \centering
    \includegraphics[width=1.1\textwidth]{LLM_en_solucion.png}
    \caption{Flujo de integración del LLM en el proceso de coordinación.}
    \label{fig:flujo_LLM}
\end{figure}

\subsection{Diagramas UML y especificaciones técnicas}

Para el diseño del sistema se consideran tres vistas principales: una vista de componentes, una vista de flujo y una vista de clases. La vista de componentes permite observar la separación entre frontend, backend, servicios internos e integraciones externas. La vista de flujo permite describir cómo un mensaje en lenguaje natural se transforma en una decisión grupal. Finalmente, la vista de clases permite representar las entidades principales utilizadas por el sistema.

Las clases principales del modelo son \texttt{Session}, \texttt{Participant}, \texttt{TimeSlot}, \texttt{TimeOption}, \texttt{AvailabilityCell}, \texttt{ChannelMessage} y \texttt{ChannelConfig}. La clase central es \texttt{Session}, ya que representa una coordinación específica. Esta contiene participantes, mensajes, disponibilidades, opciones calculadas, matriz de disponibilidad y, en caso de confirmación, una opción seleccionada y un resumen final de la decisión.

En este diseño, \texttt{Participant} solo corresponde a una persona detectada en una sesión, no a una cuenta global. Si se requiere en un futuro deberia crearse la clase \texttt{User} para soportar esto.

A nivel técnico, el MVP utiliza React y TypeScript para el frontend, FastAPI y Python para el backend, comunicación mediante HTTP/REST, persistencia local en JSON, extracción mediante LLM configurable y un motor determinístico para el cálculo de opciones. Además, se contempla una integración mínima con un canal de comunicación externo y una integración con calendario mediante generación de archivo \texttt{.ics}.

\begin{figure}[htbp]
    \centering
    \includegraphics[width=1.1\textwidth]{Clases.png}
    \caption{Diagrama de clases principales del sistema.}
    \label{fig:clases}
\end{figure}
